\documentclass[english, parskip=full]{scrartcl}
\usepackage[utf8]{inputenc}  % Kodierung der Datei
\usepackage[english]{babel}  % Sprache auf Deutsch 
\usepackage{color}
\usepackage{graphicx, gensymb}
\usepackage{amsfonts, cancel, amssymb}
\usepackage{amsmath, amsthm, array, esint}
\usepackage{booktabs, multirow}  % schönere Tabellen
\usepackage{caption}  % Anpassung der Captions
\usepackage{scrlayer-scrpage}    % Manipulation des Seitenstils
%\pagestyle{scrheadings}  % Stil für die Seite setzen
%\automark{section}  % Abschnittsnamen als Seitenbeschriftung 
%\setheadsepline{.5pt}    % Kopzeile durch Linie abtrennen
\usepackage[hidelinks]{hyperref}  % Links und weitere PDF-Features
\usepackage{typearea, tikz, pgffor, verbatim}
\usetikzlibrary{calc, quotes}
\usepackage[cache=false, outputdir=build]{minted}
\usepackage{placeins} % provides \FloatBarrier

\definecolor{bg}{rgb}{0.9,0.9,0.9}
\newcommand\numberthis{\addtocounter{equation}{1}\tag{\theequation}}
\newcommand{\bra}{}
\newcommand{\kom}{}
\newcommand{\brak}{}
\newcommand{\braket}{}
\newcommand{\intx}{}
\newcommand{\intr}{}
\newcommand{\kla}{}
\newcommand{\klam}{}
\newcommand{\klammer}{}
\newcommand{\oper}{}
\newcommand{\tecxt}{}
\newcommand{\Bra}{}
\newcommand{\Ket}{}
\renewcommand{\i}{\text{i}}
\newcommand{\e}{\text{e}}
\newcommand*{\Fourier}{\textsc{Fourier}\ }
\newcommand*{\F}{\mathcal{F}}
\newcommand*{\sinc}{\text{sinc\,}}
\newcommand{\conv}{\mathop{\scalebox{3.5}{\raisebox{-0.38ex}{\makebox[0.5\width][c]{$\ast$}}}}\limits}

\newcommand*{\titel}{04 Gauss-Elimination}
\newcommand*{\autor}{Joachim Schwardt und Julian Fleck}

\hypersetup{pdfauthor={\autor}, pdftitle={\titel}}

%\titlehead{titlehead}
\subject{Numerik I - Aufgaben}
\title{\titel}
%\author{\autor}
\author{\begin{tabular}{l|l}
Joachim Schwardt & 4768711 \\ 
Julian Fleck & 4759587\\ 
\end{tabular}}
%\author{\begin{tabular}{|l|l|}
%\hline
%Joachim Schwardt & 4768711 \\ \hline
%Julian Fleck & 4759587\\ \hline
%\end{tabular}}
\date{\today}

\begin{document}

%\begin{titlepage}
\maketitle  % Titel setzen
%\tableofcontents  % Inhaltsverzeichnis setzen
%\end{titlepage}

\section*{Aufgabe 1}
Gegeben ist die Matrix
\begin{align}
A &= \begin{pmatrix}
2 & -1 & 0 & 0 \\
-1 & 2 & -1 & 0 \\
0 & -1 & 2 & -1 \\
0 & 0 & -1 & 2
\end{pmatrix} = LU,
\end{align}
die gemäß LU-Faktorisierung in eine untere Dreiecksmatrix $L$ und eine obere Dreiecksmatrix $U$ zerlegt werden soll.
Es gilt
\begin{align}
L^{-1} &= L_3L_2L_1,
\end{align}
wobei $L_k$ eine Einheitsmatrixmatrix ist, die in der $k$-ten Spalte unterhalb der 1 die Einträge $-l_{jk}$ für $j>k$ hat.
Dabei ist $l_{jk} = \frac{a_{jk}^{(k)}}{a_{kk}^{(k)}}$.
Zunächst kann man $L_1$ ablesen, da
\begin{align}
L_1 = \begin{pmatrix}
1 & 0 & 0 & 0 \\
-l_{21} & 1 & 0 & 0 \\
-l_{31} & 0 & 1 & 0 \\
-l_{41} & 0 & 0 & 1 \\
\end{pmatrix} = \begin{pmatrix}
1 & 0 & 0 & 0 \\
1/2 & 1 & 0 & 0 \\
0 & 0 & 1 & 0 \\
0 & 0 & 0 & 1 \\
\end{pmatrix}.
\end{align}
Dann ist $A^{(2)}$ gegeben durch $L_1A^{(1)} = L_1A$, also
\begin{align}
A^{(2)} &=  \begin{pmatrix}
1 & 0 & 0 & 0 \\
1/2 & 1 & 0 & 0 \\
0 & 0 & 1 & 0 \\
0 & 0 & 0 & 1 \\
\end{pmatrix}\begin{pmatrix}
2 & -1 & 0 & 0 \\
-1 & 2 & -1 & 0 \\
0 & -1 & 2 & -1 \\
0 & 0 & -1 & 2
\end{pmatrix} = \begin{pmatrix}
2 & -1 & 0 & 0 \\
0 & 3/2 & -1 & 0 \\
0 & -1 & 2 & -1 \\
0 & 0 & -1 & 2
\end{pmatrix}.
\end{align}
Rekursiv folgen dann
\begin{align}
L_2 &= \begin{pmatrix}
1 & 0 & 0 & 0 \\
0 & 1 & 0 & 0 \\
0 & -l_{32} & 1 & 0 \\
0 & -l_{42} & 0 & 1 \\
\end{pmatrix} = \begin{pmatrix}
1 & 0 & 0 & 0 \\
0 & 1 & 0 & 0 \\
0 & 2/3 & 1 & 0 \\
0 & 0 & 0 & 1 \\
\end{pmatrix},\\
A^{(3)} &= L_2A^{(2)} = \begin{pmatrix}
1 & 0 & 0 & 0 \\
0 & 1 & 0 & 0 \\
0 & 2/3 & 1 & 0 \\
0 & 0 & 0 & 1 \\
\end{pmatrix}\begin{pmatrix}
2 & -1 & 0 & 0 \\
0 & 3/2 & -1 & 0 \\
0 & -1 & 2 & -1 \\
0 & 0 & -1 & 2
\end{pmatrix} = \begin{pmatrix}
2 & -1 & 0 & 0 \\
0 & 3/2 & -1 & 0 \\
0 & 0 & 4/3 & -1 \\
0 & 0 & -1 & 2
\end{pmatrix},\\
L_3 &= \begin{pmatrix}
1 & 0 & 0 & 0 \\
0 & 1 & 0 & 0 \\
0 & 0 & 1 & 0 \\
0 & 0 & -l_{43} & 1 \\
\end{pmatrix} = \begin{pmatrix}
1 & 0 & 0 & 0 \\
0 & 1 & 0 & 0 \\
0 & 0 & 1 & 0 \\
0 & 0 & 3/4 & 1 \\
\end{pmatrix},\\
A^{(4)} &= L_3A^{(3)} = \begin{pmatrix}
1 & 0 & 0 & 0 \\
0 & 1 & 0 & 0 \\
0 & 0 & 1 & 0 \\
0 & 0 & 3/4 & 1 \\
\end{pmatrix}\begin{pmatrix}
2 & -1 & 0 & 0 \\
0 & 3/2 & -1 & 0 \\
0 & 0 & 4/3 & -1 \\
0 & 0 & -1 & 2
\end{pmatrix} = \begin{pmatrix}
2 & -1 & 0 & 0 \\
0 & 3/2 & -1 & 0 \\
0 & 0 & 4/3 & -1 \\
0 & 0 & 0 & 5/4
\end{pmatrix}.
\end{align}
Dann gelten
%Dann gelten (das stimmt nicht, siehe unten!)
%\begin{align}
%L^{-1} = \begin{pmatrix}
%1 & 0 & 0 & 0 \\
%-l_{21} & 1 & 0 & 0 \\
%-l_{31} & -l_{32} & 1 & 0 \\
%-l_{41} & -l_{42} & -l_{43} & 1 \\
%\end{pmatrix} = \begin{pmatrix}
%1 & 0 & 0 & 0 \\
%1/2 & 1 & 0 & 0 \\
%0 & 2/3 & 1 & 0 \\
%0 & 0 & 3/4 & 1 \\
%\end{pmatrix}.
%\end{align}
%und $U = A^{(4)}$.
%Tatsächlich ist 
\begin{align}
L &= \begin{pmatrix}
1 & 0 & 0 & 0 \\
l_{21} & 1 & 0 & 0 \\
l_{31} & l_{32} & 1 & 0 \\
l_{41} & l_{42} & l_{43} & 1 \\
\end{pmatrix} = \begin{pmatrix}
1 & 0 & 0 & 0 \\
-1/2 & 1 & 0 & 0 \\
0 & -2/3 & 1 & 0 \\
0 & 0 & -3/4 & 1 \\
\end{pmatrix},
\end{align}
und $U = A^{(4)}$.
Überprüfen des Ergebnis gibt
\begin{align}
LU &= \begin{pmatrix}
1 & 0 & 0 & 0 \\
-1/2 & 1 & 0 & 0 \\
0 & -2/3 & 1 & 0 \\
0 & 0 & -3/4 & 1 \\
\end{pmatrix}\begin{pmatrix}
2 & -1 & 0 & 0 \\
0 & 3/2 & -1 & 0 \\
0 & 0 & 4/3 & -1 \\
0 & 0 & 0 & 5/4
\end{pmatrix} = \\
&= \begin{pmatrix}
2 & -1 & 0 & 0 \\
-1 & 2 & -1 & 0 \\
0 & -1 & 2 & -1 \\
0 & 0 & -1 & 2
\end{pmatrix} = A,
\end{align}
womit $L$ und $U$ also korrekt bestimmt sind.
\\
Die Inverse von $L$ ist
\begin{align}
L^{-1} = L_3L_2L_1 = \begin{pmatrix}
1 & 0 & 0 & 0 \\
1/2 & 1 & 0 & 0 \\
1/3 & 2/3 & 1 & 0 \\
1/4 & 1/2 & 3/4 & 1 \\
\end{pmatrix} \neq L_\tecxt\text{oben}^{-1}.
\end{align}
Für die Inverse kann man also scheinbar nicht einfach die Vorzeichen unterhalb der Diagonalen drehen, das funktioniert nur für $L_k$ und $L_k^{-1}$.
Für das Berechnen von $U^{-1}$ kann man Rückwärts-Substitution verwenden.
Die Diagonaleinträge sind zunächst einfach der Kehrwert der Diagonaleinträge von $U$.
Dann löst man iterativ
\begin{align}
U^{-1} &= \begin{pmatrix}
1/2 & u_{12} & u_{13} & u_{14} \\
0 & 2/3 & u_{23} & u_{24} \\
0 & 0 & 3/4 & u_{34}  \\
0 & 0 & 0 & 4/5
\end{pmatrix}\begin{pmatrix}
2 & -1 & 0 & 0 \\
0 & 3/2 & -1 & 0 \\
0 & 0 & 4/3 & -1 \\
0 & 0 & 0 & 5/4
\end{pmatrix} \overset{!}{=} \begin{pmatrix}
1 & 0 & 0 & 0 \\ 
0 & 1 & 0 & 0 \\
0 & 0 & 1 & 0 \\ 
0 & 0 & 0 & 1 \\
\end{pmatrix}.
\end{align}
es ist $-\frac{3}{4} + \frac{5}{4} u_{34} = 0$, also $u_{34} = \frac{3}{5}$.
Analog findet man $u_{23} = \frac{2}{4} = \frac{1}{2}$ und $u_{12} = \frac{1}{3}$. 
Es ist $-u_{23} + \frac{5}{4} u_{24} = 0$, also $u_{24} = \frac{2}{5}$.
Analog $u_{13} = \frac{1}{4}$.
Zuletzt ist noch $-u_{13} + \frac{5}{4}u_{14} = 0$, also $u_{14} = \frac{1}{5}$.
Damit können wir die Inverse von $A$ einfach bestimmen, da
\begin{align}
A^{-1} &= U^{-1}L^{-1} = \begin{pmatrix}
1/2 & 1/3 & 1/4 & 1/5 \\
0 & 2/3 & 2/4 & 2/5 \\
0 & 0 & 3/4 & 3/5  \\
0 & 0 & 0 & 4/5
\end{pmatrix}\begin{pmatrix}
1 & 0 & 0 & 0 \\
1/2 & 1 & 0 & 0 \\
1/3 & 2/3 & 1 & 0 \\
1/4 & 1/2 & 3/4 & 1 \\
\end{pmatrix} = \\
&= \begin{pmatrix}
0.8 & 0.6 & 0.4 & 0.2 \\
0.6 & 1.2 & 0.8 & 0.4 \\
0.4 & 0.8 & 1.2 & 0.6 \\
0.2 & 0.4 & 0.6 & 0.8 \\
\end{pmatrix}.
\end{align}
Die Einträge liegen alle in einer ähnlichen Größenordnung, es sollten also keine numerischen Problem auftreten, wenn man $A^{-1}$ zum Lösen eines Gleichungssystems $Ax=b$ verwendet.
%Für andere Matrizen sind Einträge der Inversen häufig extrem klein oder groß, was schnell die Rechengenauigkeit erschöpfen kann.
%Beispielsweise hat die Matrix

% B = (np.linspace(1, n**2, n**2)**2).reshape((n,n))

% n = 4
% A = np.diag(np.full(n, 2.0)) + np.diag(np.full(n-1, -1), k=1) + np.diag(np.full(n-1, -1), k=-1)
% L = np.eye(n) + np.diag(np.full(n-1, [1/k - 1 for k in range(2, n+1, 1)]), k=-1)
% U = np.diag(np.full(n, [1/k + 1 for k in range(1, n+1, 1)])) + np.diag(np.full(n-1, -1), k=1)

%for row in range(n):
%    print(" & ".join(f"{Ainv[row, col]:.1f}" for col in range(n)) + r" \\")


\newpage
\section*{Aufgabe 2}
Gegeben ist ein ``Computer'' mit einer Rechengenauigkeit von vier signifikanten Stellen im Dezimalsystem.
\subsection*{a)}
Zu lösen ist das Gleichungssystem
\begin{align}
1,000\cdot 10^{-3} x_1 + 2,200\cdot 10^1 x_2 &= 4,000\cdot 10^1 \\
4,420\cdot 10^2 x_1 + 4,000\cdot 10^0 x_2 &= 1,200\cdot 10^2.
\end{align}
Die numerische Lösung mittel Python liefert $x_1 = 2,5504\cdot 10^{-1}$ und $x_2 = 1,8182\cdot 10^0$ (gerundet auf fünf signifikante Stellen).
\\
Ohne Pivotisierung ist $\ell_1 = -\frac{a_{21}}{a_{11}} = -\frac{4,420\cdot 10^2}{1,000\cdot 10^{-3}} = -4,420\cdot 10^5$:
\begin{align*}
\tecxt\text{(II)} - \ell_1\cdot \tecxt\text{(I)   :} \ \ \  &\kla\left( 4,420\cdot 10^2 - 4,420\cdot 10^5 \cdot 1,000\cdot 10^{-3} \right)x_1 + \\
& + \kla\left( 4,000\cdot 10^0 -  4,420\cdot 10^5\cdot 2,200\cdot 10^1 \right)x_2 = \\
&= 1,200\cdot 10^2 - 4,420\cdot 10^5\cdot 4,000\cdot 10^1 \numberthis\\
& (0,000\cdot 10^2)x_1 + \kla\left( 4,000\cdot 10^0 - 9,724\cdot 10^6 \right)x_2 = -1,768\cdot 10^7 \numberthis \\
& -9,724\cdot 10^6 x_2 = -1,768\cdot 10^7 \numberthis \\
& x_2 = \frac{1,768\cdot 10^7}{9,724\cdot 10^6} = 1,818 \cdot 10^0. \numberthis
\end{align*}
Für $x_1$ ergibt sich dann
\begin{align}
x_1 &= \frac{1}{1,000\cdot 10^{-3}}\kla\left( 4,000\cdot 10^1 - 2,200\cdot 10^1 \cdot 1,818\cdot 10^0 \right) = \\
&= \frac{1}{1,000\cdot 10^{-3}}\kla\left( 4,000\cdot 10^1 - 4,000\cdot 10^1 \right) = \\
&=\frac{0,000\cdot 10^1}{1,000\cdot 10^{-3}} = 0,000\cdot 10^4.
\end{align}
Dieses Ergebnis ist im Rahmen der Rechengenauigkeit also ``falsch'', da die Abweichung nicht in der letzten Stelle auftritt (anders ausgedrückt: $x_1=0$ ist völlig daneben :) ). 
\\
Mit Spaltenpivotisierung vertauschen wir die beiden Zeilen, da $|a_{11}| < |a_{21}|$.
Dann ist $\ell_1 = -\frac{a_{21}}{a_{11}} = -\frac{1,000\cdot 10^{-3}}{4,420\cdot 10^2}  = -2,262\cdot 10^{-6}$:
\begin{align*}
\tecxt\text{(II)} - \ell_1\cdot \tecxt\text{(I)   :} \ \ \  &\kla\left( 1,000\cdot 10^{-3} - 2,262\cdot 10^{-6} \cdot 4,420\cdot 10^2 \right)x_1 + \\
& + \kla\left( 2,200\cdot 10^1 - 2,262\cdot 10^{-6} \cdot 4,000\cdot 10^0 \right)x_2 = \\
&= 4,000\cdot 10^1 - 2.262\cdot 10^{-6}\cdot 1,200\cdot 10^2 \numberthis\\
& (1,000\cdot 10^{-3} - 9,998\cdot 10^{-4})x_1 + \\
& + \kla\left( 2,200\cdot 10^1 - 9,048\cdot 10^{-6} \right)x_2 = 4,000\cdot 10^1 - 2.714\cdot 10^{-4} \numberthis\\
& 2,200\cdot 10^1 x_2 = 4,000\cdot 10^1  \numberthis\\
x_2 &= \frac{4,000\cdot 10^1}{2,200\cdot 10^1} = 1,818\cdot 10^0. \numberthis
\end{align*}
Für $x_1$ ergibt sich dann
\begin{align}
x_1 &= \frac{1}{4,420\cdot 10^{2}}\kla\left( 1,200\cdot 10^2 - 4,000\cdot 10^0 \cdot 1.818\cdot 10^0 \right) = \\
&= \frac{1}{4,420\cdot 10^{2}}\kla\left( 1,200\cdot 10^2 - 7,272\cdot 10^0 \right) = \\
&=\frac{1,127\cdot 10^2}{4,420\cdot 10^{2}} = 2,550\cdot 10^{-1}.
\end{align}
Dieses Ergebnis ist im Rahmen der Rechengenauigkeit also ``richtig''.
\\
Bei vollständiger Pivotisierung ändert sich nichts, da $4,420\cdot 10^2$ bereits der größere Faktor in der -- nach der Spaltenpivotisierung nun ``ersten'' --  Zeile ist.
\subsection*{b)}
Wir multiplizieren die erste Zeile des Gleichungssystems aus a) mit $1,000\cdot 10^6$ und die zweite mit $2,500\cdot 10^{-1}$.
Zu lösen ist also das Gleichungssystem
\begin{align}
1,000\cdot 10^{3} x_1 + 2,200\cdot 10^7 x_2 &= 4,000\cdot 10^7 \\
1,105\cdot 10^2 x_1 + 1,000\cdot 10^0 x_2 &= 3,000\cdot 10^1.
\end{align}
Ohne Pivotisierung ist $\ell_1 = -\frac{a_{21}}{a_{11}} = -\frac{1,105\cdot 10^2}{1,000\cdot 10^{3}} = -1,105\cdot 10^{-1}$:
\begin{align*}
\tecxt\text{(II)} - \ell_1\cdot \tecxt\text{(I)   :} \ \ \  &\kla\left( 1,105\cdot 10^2 - 1,105\cdot 10^{-1} \cdot 1,000\cdot 10^{3} \right)x_1 + \\
& + \kla\left( 1,000\cdot 10^0 -  1,105\cdot 10^{-1}\cdot 2,200\cdot 10^7 \right)x_2 = \\
&= 3,000\cdot 10^1 - 1,105\cdot 10^{-1}\cdot 4,000\cdot 10^7 \numberthis\\
& \kla\left( 1,000\cdot 10^0 - 2,431\cdot 10^6 \right)x_2 = -4,420\cdot 10^6 \numberthis \\
&  - 2,431\cdot 10^6 x_2 = -4,420\cdot 10^6  \numberthis \\
& x_2 = \frac{4,420\cdot 10^6}{2,431\cdot 10^6} = 1.818 \cdot 10^0. \numberthis
\end{align*}
Für $x_1$ ergibt sich dann
\begin{align}
x_1 &= \frac{1}{1,000\cdot 10^{3}}\kla\left( 4,000\cdot 10^7 - 2,200\cdot 10^7 \cdot 1,818\cdot 10^0 \right) = \\
&= \frac{1}{1,000\cdot 10^{3}}\kla\left( 4,000\cdot 10^7 - 4,000\cdot 10^7 \right) = \\
&=\frac{0,000\cdot 10^7}{1,000\cdot 10^{3}} = 0,000\cdot 10^4.
\end{align}
Dieses Ergebnis ist im Rahmen der Rechengenauigkeit also ``falsch''. 
\\
Mit Spaltenpivotisierung ändert sich nichts, da bereits $|a_{11}| > |a_{21}|$.
Bei vollständiger Pivotisierung vertauschen wir die Spalten, da $|a_{12}| > |a_{11}|$.
Im Endergebnis müssen wir zudem noch $x_1$ und $x_2$ tauschen.
Hier ist dann $\ell_1 = -\frac{a_{21}}{a_{11}} = -\frac{1,000\cdot 10^0}{2,200\cdot 10^{7}}  = -4.545\cdot 10^{-8}$:
\begin{align*}
\tecxt\text{(II)} - \ell_1\cdot \tecxt\text{(I)   :} \ \ \  &\kla\left( 1,000\cdot 10^{0} - 4.545\cdot 10^{-8} \cdot 2,200\cdot 10^{7} \right)x_1 + \\
& + \kla\left( 1,105\cdot 10^2 - 4.545\cdot 10^{-8} \cdot 1,000\cdot 10^3 \right)x_2 = \\
&= 3,000\cdot 10^1 - 4.545\cdot 10^{-8}\cdot 4,000\cdot 10^7 \numberthis\\
& (1,000\cdot 10^{0} - 9,999\cdot 10^{-1})x_1 + \\
& + \kla\left( 1,105\cdot 10^2 - 4,545\cdot 10^{-4} \right)x_2 = 3,000\cdot 10^1 - 1,818\cdot 10^{0} \numberthis\\
& 1,105\cdot 10^2 x_2 = 2,818\cdot 10^1  \numberthis\\
x_2 &= \frac{2,818\cdot 10^1}{1,105\cdot 10^2} = 2,550\cdot 10^{-1}. \numberthis
\end{align*}
Für $x_1$ ergibt sich dann
\begin{align}
x_1 &= \frac{1}{2,200\cdot 10^{7}}\kla\left( 4,000\cdot 10^7 - 1,000\cdot 10^3 \cdot 2,550\cdot 10^{-1} \right) = \\
&= \frac{1}{2,200\cdot 10^{7}}\kla\left( 4,000\cdot 10^7 - 2,550\cdot 10^{2} \right) = \\
&=\frac{4,000\cdot 10^7}{2,200\cdot 10^{7}} = 1,818\cdot 10^{0}.
\end{align}
Dieses Ergebnis ist im Rahmen der Rechengenauigkeit also wieder ``richtig''.
\subsection*{c)}
Für a) mit Spaltenpivotisierung
\begin{align}
P_1 &= \begin{pmatrix}
0 & 1 \\ 1 & 0
\end{pmatrix}
\end{align}
Für a) mit vollständiger Pivotisierung
\begin{align}
P_1 &= \begin{pmatrix}
0 & 1 \\ 1 & 0
\end{pmatrix}&&&P_2 &= \begin{pmatrix}
1 & 0 \\ 0 & 1
\end{pmatrix}
\end{align}
%Exakte Lösung:
%\begin{align}
%\begin{pmatrix}
%\frac{1}{1000} & 22 \\
%442 & 4
%\end{pmatrix}\begin{pmatrix}
%x_1 \\ x_2
%\end{pmatrix} &= \begin{pmatrix}
%40 \\ 120
%\end{pmatrix}\\
%\begin{pmatrix}
%\frac{1}{1000} & 22 \\
%0 & 4 - 22\cdot 1000\cdot 442
%\end{pmatrix}\begin{pmatrix}
%x_1 \\ x_2
%\end{pmatrix} &= \begin{pmatrix}
%40 \\ 120 - 40\cdot 1000\cdot 442
%\end{pmatrix}
%\end{align}
%$x$
%Also ist $x_2 = \frac{120 - 1768\cdot 10000}{4 - 9724\cdot 1000} = 1,81817$

\newpage
\section*{Aufgabe 3}
Sei $A\in\mathbb{R}^{n\times n}$ invertierbar und $\|\cdot\|$ eine induzierte Matrixnorm. 
Die Kondition von $A$ sei definiert als
\begin{align}
\kappa(A) &:= \|A\|\|A^{-1}\|.
\end{align}
\subsection*{a)}
Die induzierte Matrixnorm ist definiert als
\begin{align}
\|A\| &:= \sup_{x\neq 0}\frac{\|Ax\|}{\|x\|}.
\end{align}
Aus der Homogenität der Norm folgt dann direkt die Darstellung
\begin{align}
\|A\| &= \sup_{x\neq 0}\left\|A\frac{x}{\|x\|}\right\| = \max_{\|y \| = 1}\|A y\|. \label{eqn:norm of A}
\end{align}
Die Norm der Inversen ist mit $y := A^{-1}x$ für beliebige $x\in\mathbb{R}^n$ gegeben durch
\begin{align}
\|A^{-1}\| &= \sup_{x\neq 0}\frac{\|A^{-1}x\|}{\| x \|} = \sup_{Ay \neq 0} \frac{\|A^{-1}Ay\|}{\|Ay\|} = \sup_{y \neq 0} \frac{\|y\|}{\|Ay\|},
\end{align}
da $Ay\neq 0$ aufgrund der Invertierbarkeit von $A$ äquivalent zu $y\neq 0$ ist.
Weiterhin können wir schreiben
\begin{align}
\|A^{-1}\| &= \sup_{y\neq 0}\frac{1}{\left\| A \frac{y}{\|y\|}\right\|} = \max_{\|x\| = 1} \frac{1}{\|Ax\|} = \kla\left( \min_{\|x\| = 1}\|Ax\| \right)^{-1}. \label{eqn:norm of A inverse}
\end{align}
Im letzten Schritt haben wir dabei verwendet, dass der Ausdruck $\frac{1}{\|Ax\|}$ genau dann maximal ist, wenn der Nenner $\|Ax\|$ minimal ist.
Insgesamt folgt aus \eqref{eqn:norm of A} und \eqref{eqn:norm of A inverse} die Darstellung
\begin{align}
\kappa (A) &= \frac{\underset{\|x\| = 1}{\max}\|Ax\|}{\underset{\|x\| = 1}{\min}\|Ax\|} \label{eqn:kappa A max min formel}
\end{align}
für die Kondition einer invertierbaren Matrix $A$.
\subsection*{b)}
Da die Norm per Definition positiv definit ist, gilt 
\begin{align}
\kappa(A) &\ge 0.
\end{align}
Insbesondere ist $\kappa (A) = 0$ äquivalent zu $\|A\| = 0$, und damit zu $A = 0$.
Aber die Nullmatrix $A=0$ ist nicht invertierbar, also muss $\kappa(A) > 0$ sein.
\\
Sei $\alpha\in\mathbb{R}\setminus\{0\}$. Dann folgt aus der Eigenschaft \eqref{eqn:kappa A max min formel} direkt
\begin{align}
\kappa (\alpha A) = \frac{\underset{\|x\| = 1}{\max}\|\alpha Ax\|}{\underset{\|x\| = 1}{\min}\|\alpha Ax\|} = \frac{\underset{\|x\| = 1}{\max}|\alpha|\| Ax\|}{\underset{\|x\| = 1}{\min}|\alpha|\| Ax\|} = \kappa (A).
\end{align}
\subsection*{c)}
Sei $A_k$ eine Folge invertierbarer Matrizen, die gegen eine singuläre Matrix $A\neq 0\in\mathbb{R}^{n\times n}$ konvergiert.
Dann existiert eine Folge von Eigenwerten der $A_k$ mit $\lambda_k \rightarrow 0$ für $k\rightarrow 0$, da aus $\det A = 0$ folgt, dass mindestens ein Eigenwert gleich Null ist.
Insbesondere kann der zugehörige Eigenvektor $x_k\in\mathbb{R}^n\setminus\{0\}$ mit $A_kx_k = \lambda_k x_k$ auf Eins normiert werden.
Sei also $\|x_k\| = 1$.
Dann gilt
\begin{align}
\frac{1}{\|A_k^{-1}\|} = \min_{\|x\| = 1} \|A_kx\| \le \|A_kx_k\| = \|\lambda_kx_k\| = |\lambda_k| \overset{k\rightarrow\infty}{\longrightarrow} 0.
\end{align}
Der entscheidende Punkt ist, dass wir $A\neq 0$ annehmen.
Da die Norm definit ist, sind gilt $\|A_k\| \neq 0$ für alle $k$, aber insbesondere auch $\|A\| \neq 0$.
Die Norm aller Folgenglieder und auch des Grenzwertes ist also streng positiv, weshalb
\begin{align}
\lim_{k\rightarrow\infty} \kappa (A_k) &= \|A\| \lim_{k\rightarrow\infty} \|A_k^{-1}\| 
\end{align}
Aus der Abschätzung für die Norm der inversen von $A_k$ folgt dann aber
\begin{align}
\lim_{k\rightarrow\infty} \kappa (A_k) &\ge \|A\| \lim_{k\rightarrow \infty}\frac{1}{|\lambda_k|} = \infty.
\end{align}
Ist $A\neq 0$ tatsächlich eine wichtige Voraussetzung? 
Betrachte dazu das Beispiel 
\begin{align}
A_k &:= \begin{pmatrix}
\frac{1}{k} & \dots & 0 \\ 
\vdots & \ddots & \vdots \\ 
0 & \dots & \frac{1}{k}
\end{pmatrix}.
\end{align}
%Offensichtlich ist $\det A_k = \frac{1}{k^n} \rightarrow 0$ für $k\rightarrow\infty$.
Für alle $k\in\mathbb{N}$ gilt
\begin{align}
\kappa (A_k) &= \frac{\underset{\|x\| = 1}{\max}\|A_kx\|}{\underset{\|x\| = 1}{\min}\|A_kx\|} = \frac{\underset{\|x\| = 1}{\max}\frac{1}{k}\|x\|}{\underset{\|x\| = 1}{\min}\frac{1}{k}\|x\|} = 1.
\end{align}
Insbesondere gilt also 
$$\forall\epsilon > 0\exists K\in\mathbb{N}\forall k\ge K: |\kappa(A_k) - 1| < \epsilon, $$
da offensichtlich $|\kappa(A_k) - 1 | = 0 < \epsilon$ für beliebige $k$ ist.
Dies ist aber gerade die Definition des Grenzwertes, also ist 
\begin{align}
\lim_{k\rightarrow\infty}\kappa (A_k) &= 1.
\end{align}
Gleichzeitig ist aber
\begin{align}
\| A_k - 0\| &= \max_{\|x\| = 1} \|A_kx\| = \frac{1}{k} \overset{k\rightarrow\infty}{\longrightarrow} 0.
\end{align}
Die Folge invertierbarer Matrizen $A_k$ konvergiert in der Norm also gegen die Nullmatrix, die natürlich singulär ist.
In diesem Fall gilt der Satz $\kappa(A_k)\rightarrow\infty$ also nicht.
Anmerkung: Letzte Woche wurde die Aufgabenstellung korrigiert, wir hatten das Gegenbeispiel zu dem Zeitpunkt aber schon aufgeschrieben -- deshalb lassen wir das jetzt einfach drin :)

\subsection*{d)}
%Allgemein gilt zunächst
%\begin{align}
%\|A \| &= \max_{\|x\| = 1} \|Ax\| \ge \|Ax_\tecxt\text{max}\| = |\lambda_\tecxt\text{max}| \label{eqn:norm A ge lambda max}
%\end{align}
%für den Eigenvektoren $x_\tecxt\text{max}$ von $A$ mit (betragsmäßig) maximalem Eigenwert $\lambda_\tecxt\text{max}$.
%Außerdem gilt analog für den normierten Eigenvektor $x_\tecxt\text{min}$ mit minimalen Eigenwert $\lambda_\tecxt\text{min}$
%\begin{align}
%\frac{1}{\| A^{-1} \|} &= \min_{\|x\|=1} \|Ax\| \le \|Ax_\tecxt\text{min}\| = |\lambda_\tecxt\text{min}|. \label{eqn:norm 1 over A inverse le lambda min}
%\end{align}
Sei nun $A$ zusätzlich symmetrisch und $\|\cdot\| = \|\cdot\|_2$.
Aus der Symmetrie folgt, dass $A$ diagonalisierbar ist.
Also existiert eine orthogonale Matrix $S\in\mathbb{R}^{n\times n}$, sodass
\begin{align}
A &= S^{-1}DS,
\end{align}
wobei $D\in\mathbb{R}^{n\times n}$ eine diagonale Matrix mit den Eigenwerten $\lambda_n$ von $A$ ist und $S^{-1} = S^\top$.
Insbesondere erhält $S$ die Norm, das heißt $\|Sx\|_2 = \|x\|_2$ für beliebige $x\in\mathbb{R}^n$, da
\begin{align}
\|Sx\|_2^2 &= x^\top S^\top Sx = x^\top x = \|x\|_2^2.
\end{align}
Damit folgt
\begin{align}
\|A\| &= \max_{\|x\| = 1} \|S^{-1}DSx\| = \max_{\|x\| = 1} \|DSx\| \max_{\|y\| = \|Sx\| = 1} \|Dy\| = |\lambda_\tecxt\text{max}|,
\end{align}
%und zusammen mit \eqref{eqn:norm A ge lambda max} die Gleichheit 
%\begin{align}
%\|A\| &= |\lambda_\tecxt\text{max}|.
%\end{align}
Analog gilt für die Inverse
%Analog gilt für die Inverse $A^{-1} = (S^{-1}DS)^{-1} = S^{-1}(S^{-1}D)^{-1} = S^{-1}D^{-1}S$ 
\begin{align}
\frac{1}{\|A^{-1}\|} &= \min_{\|x\| = 1} \|S^{-1}DSx\| = \min_{\|x\| = 1} \|DSx\| = \min_{\|y\| = \|Sx\| = 1} \|Dy\| = |\lambda_\tecxt\text{min}|.
\end{align}
Insgesamt gilt also
\begin{align}
\kappa (A) &= \|A\| \|A^{-1}\| = |\lambda_\tecxt\text{max}|\frac{1}{|\lambda_\tecxt\text{min}|} = \left\vert\frac{\lambda_\tecxt\text{max}}{\lambda_\tecxt\text{min}}\right\vert.
\end{align}

%Es folgt
%\begin{align}
%\|A\| &= \|S^{-1}DS\| = \sup_{x\neq 0}\frac{\|S^{-1}DSx\|}{\|x\|} = \sup_{y = Sx \neq 0} \frac{\|S^{-1}Dy\|}{\|S^{-1}y\|}
%\end{align}
%\begin{align}
%\|A\| &= \max_{\|x\|=1} \|S^{-1}DSx\| = \max_{\|y\| = \|Sx\| = 1} \| S^{-1} Dy\|
%\end{align}
%\begin{align}
%\kappa (A) &= \frac{\underset{\|x\| = 1}{\max}\|Ax\|}{\underset{\|x\| = 1}{\min}\|Ax\|} = \frac{\underset{\|x\| = 1}{\max}\|S^{-1}DSx\|}{\underset{\|x\| = 1}{\min}\|S^{-1}DSx\|} = \frac{\underset{\|x\| = 1}{\max}\|Ax\|}{\underset{\|x\| = 1}{\min}\|Ax\|} = 
%\end{align}
%\begin{align}
%\|S^{-1}\| &= \sup_{x\neq 0}\frac{\|S^{-1}x\|}{\|x\|} \ge \frac{\|y\|}{\|Sy\|} &&\Rightarrow & \|S^{-1}y\|&\ge \frac{1}{\|S\|}
%\end{align}


\newpage
\section*{Aufgabe 4}
Für die Matrizen $S^n$ mit
\begin{align}
S_{jk}^n := \left\{\begin{matrix}
2 & : & j=k, \\
-1 & : & |j-k| = 1,\\
0 & : & \tecxt\text{sonst}
\end{matrix} \right. 
\end{align}
soll die Durchführbarkeit des Gauss'schen Eliminierungsverfahrens bewiesen werden.
Wir wollen dazu explizit die Einträge der oberen Dreiecksmatrix $T^n$ konstruieren, welche sich bei dem Verfahren ergibt. 
Da oberhalb der Hauptdiagonalen von $S^n$ hier alle Einträge unverändert bleibt, ist die allgemeine obere Dreiecksform von $T^n$ bekannt, wenn wir die Hauptdiagonalelemente angeben.
Diese sind $T_{jj}^n = 1 + \frac{1}{j}$ für $j\le n$.
Für $n=1$ ist offensichtlich $T_{jj}^1 = 1 + \frac{1}{j} = 2$ für $j\le n$.
Sei also nun $T_{jj}^n = 1 + \frac{1}{j}$ für $j\le n$.
Dann erhalten wir im vorletzten Schritt des Verfahrens also die Matrix
\begin{align}
\begin{pmatrix}
\ddots & \vdots & \vdots & \vdots \\
\dots & 1 + \frac{1}{n-2} & -1 & 0 \\
\dots & 0 & 1 + \frac{1}{n-1} & -1 \\
\dots & 0 & -1 & 2
\end{pmatrix}.
\end{align}
Sei nun für den letzten Schritt
\begin{align}
\ell &:= \frac{-1}{1 + \frac{1}{n-1}} = \frac{1}{n} - 1.
\end{align}
Wir subtrahieren nun $\ell$ mal die vorletzte, also $(n-1)$-te, Zeile von der letzten Zeile und erhalten letztlich
\begin{align}
T^n &= \begin{pmatrix}
\ddots & \vdots & \vdots & \vdots \\
\dots & 1 + \frac{1}{n-2} & -1 & 0 \\
\dots & 0 & 1 + \frac{1}{n-1} & -1 \\
\dots & 0 & 0 & 2 - \ell \cdot (-1)
\end{pmatrix} = \begin{pmatrix}
\ddots & \vdots & \vdots & \vdots \\
\dots & 1 + \frac{1}{n-2} & -1 & 0 \\
\dots & 0 & 1 + \frac{1}{n-1} & -1 \\
\dots & 0 & 0 & 1 + \frac{1}{n}
\end{pmatrix}.
\end{align}
Damit ist der Induktionsbeweis vollständig.


%\begin{thebibliography}{9}
%\bibitem{example} description, \url{https://www.exmapleurl}, day.\,month~2021.
%\end{thebibliography}

\end{document}
